% Template for PLoS
% Version 3.8 Apr 2026
%
% % % % % % % % % % % % % % % % % % % % % % %
%
% -- IMPORTANT NOTE
%
% This template contains comments intended 
% to minimize problems and delays during our production 
% process. Please follow the template instructions
% whenever possible.
%
% % % % % % % % % % % % % % % % % % % % % % % 
%
% Once your paper is accepted for publication, 
% PLEASE REMOVE ALL TRACKED CHANGES in this file 
% and leave only the final text of your manuscript. 
% PLOS recommends the use of latexdiff to track changes during review, as this will help to maintain a clean tex file.
% Visit https://www.ctan.org/pkg/latexdiff?lang=en for info.
%
%
% IMPORTANT
% Below are a few tips to help format manuscript according to our specifications. For more tips to help reduce the possibility of formatting errors during conversion, please see our LaTeX guidelines at http://journals.plos.org/plosone/s/latex
%
% There are no restrictions on package use within the LaTeX files except that no packages listed in the template may be deleted.
%
% Please do not include colors or graphics in the text.
%
% The manuscript LaTeX source should be contained within a single file (do not use \input, \externaldocument, or similar commands).
%
% % % % % % % % % % % % % % % % % % % % % % %
%
% -- FIGURES AND TABLES
%
% Please include tables/figure captions directly after the paragraph where they are first cited in the text.
%
% DO NOT INCLUDE GRAPHICS IN YOUR MANUSCRIPT
% - Figures should be uploaded separately from your manuscript file. 
% - Figures generated using LaTeX should be extracted and removed from the PDF before submission. 
% - Figures containing multiple panels/subfigures must be combined into one image file before submission.
% For figure citations, please use "Fig" instead of "Figure".
% See http://journals.plos.org/plosone/s/figures for PLOS figure guidelines.
%
% Tables should be cell-based and may not contain:
% - spacing/line breaks within cells to alter layout or alignment
% - do not nest tabular environments (no tabular environments within tabular environments)
% - no graphics or colored text (cell background color/shading OK)
% See http://journals.plos.org/plosone/s/tables for table guidelines.
%
% For tables that exceed the width of the text column, use the adjustwidth environment as illustrated in the example table in text below.
%
% % % % % % % % % % % % % % % % % % % % % % % %
%
% -- EQUATIONS, MATH SYMBOLS, OTHER SPECIAL FORMATTING
%
% - For bold symbols, use the \boldsymbol command and not \mathbf. 
% - For inline equations, please be sure to include all portions of an equation in the math environment.  For example, x$^2$ is incorrect; this should be formatted as $x^2$ (or $\mathrm{x}^2$ if the romanized font is desired).
% - Do not include text that is not math in the math environment. For example, CO2 should be written as CO\textsubscript{2} instead of CO$_2$.
% - Avoid capturing simple text as equations, for example $<$.
% - For inline equations, please do not include punctuation (commas, etc) within the math environment unless this is part of the equation.
% - When adding superscript or subscripts outside of brackets/braces, please group using {}.  For example, change "[U(D,E,\gamma)]^2" to "{[U(D,E,\gamma)]}^2". 
% - Do not use \cal for caligraphic font.  Instead, use \mathcal{}.
% - Avoid splitting for fragmenting a single equation into multiple objects/environments, for example $f(x)$ = $g(x)$ + $h(x)$.
% - Use math environment for all equations, do not mimic using text.
% - Avoid using single letters in math mode where possible.
% - Always use matching parentheses and delimiters, avoid mixing math and text brackets.
% - Insert spaces around inline equations to ensure proper display after conversion.
% - Do not use forced numbers for display equation labeling or suppress numbering with \nonumber. 
%
% % % % % % % % % % % % % % % % % % % % % % % % 
%
% Please contact customercare@plos.org with any questions.
%
% % % % % % % % % % % % % % % % % % % % % % % %

\documentclass[10pt,letterpaper]{article}
\usepackage[top=0.85in,left=2.75in,footskip=0.75in]{geometry}

% amsmath and amssymb packages, useful for mathematical formulas and symbols
\usepackage{amsmath,amssymb}

% Use adjustwidth environment to exceed column width (see example table in text)
\usepackage{changepage}

% Report PDF: keep figures where cited (no float to end of document)
\usepackage{float}

% textcomp package and marvosym package for additional characters
\usepackage{textcomp,marvosym}

% cite package, to clean up citations in the main text. Do not remove.
\usepackage{cite}

% Use nameref to cite supporting information files (see Supporting Information section for more info)
\usepackage{nameref,hyperref}

% line numbers
\usepackage[right]{lineno}

% ligatures disabled
\usepackage[nopatch=eqnum]{microtype}
\DisableLigatures[f]{encoding = *, family = * }

% color can be used to apply background shading to table cells only
\usepackage[table]{xcolor}

% array package and thick rules for tables
\usepackage{array}

% create "+" rule type for thick vertical lines
\newcolumntype{+}{!{\vrule width 2pt}}

% create \thickcline for thick horizontal lines of variable length
\newlength\savedwidth
\newcommand\thickcline[1]{%
  \noalign{\global\savedwidth\arrayrulewidth\global\arrayrulewidth 2pt}%
  \cline{#1}%
  \noalign{\vskip\arrayrulewidth}%
  \noalign{\global\arrayrulewidth\savedwidth}%
}

% \thickhline command for thick horizontal lines that span the table
\newcommand\thickhline{\noalign{\global\savedwidth\arrayrulewidth\global\arrayrulewidth 2pt}%
\hline
\noalign{\global\arrayrulewidth\savedwidth}}


% Remove comment for double spacing
%\usepackage{setspace} 
%\doublespacing

% Text layout
\usepackage{ragged2e}
\setlength{\parindent}{0.5cm}
\textwidth 5.25in
\textheight 8.75in

% Bold the 'Figure #' in the caption and separate it from the title/caption with a period
% Captions will be left justified
\usepackage[aboveskip=1pt,font=normal,labelfont=bf,labelsep=period,justification=raggedright,singlelinecheck=off]{caption}
\renewcommand{\figurename}{Fig}

% Please use the included `plos2025.bst` as your BibTeX style
\bibliographystyle{assets/plos2025}

% Remove brackets from numbering in List of References
\makeatletter
\renewcommand{\@biblabel}[1]{\quad#1.}
\makeatother

% QSLP-style bold paragraph starter (see report_outline.md)
\newcommand{\methodpara}[1]{\par\vspace{0.85em}\noindent\textbf{#1}\ }

% Appendix cross-ref: subsection labels render as A.1, B.2, etc. after \appendix
\newcommand{\apdx}[1]{Appendix~\ref{#1}}
\newcommand{\apptabdesc}[1]{\par\vspace{0.45em}\noindent#1\par\vspace{0.35em}}
\newcommand{\apptabular}[1]{%
  \par\vspace{0.35em}%
  \begin{center}%
  #1%
  \end{center}%
  \par\vspace{0.5em}%
}
\newcommand{\appfigdesc}[1]{\par\smallskip\noindent#1\par\smallskip}



% Header and Footer with logo
\usepackage{lastpage,fancyhdr,graphicx}
\usepackage{epstopdf}

% Local report PDF with embedded figures (not for PLOS upload):
%   bash build/compile_report.sh  →  plos_my_paper_report.pdf
% Submission build (caption-only):  pdflatex plos_my_paper.tex
\makeatletter
\@ifundefined{draftfigureson}{%
  \newif\ifdraftfigures
  \draftfiguresfalse
}{%
  \newif\ifdraftfigures
  \draftfigurestrue
}%
\@ifundefined{reportbuild}{%
  \newif\ifreportmode
  \reportmodefalse
}{%
  \newif\ifreportmode
  \reportmodetrue
}%
\makeatother
\ifreportmode
  \usepackage{titlesec}
  \setcounter{secnumdepth}{2}
  \titleformat{\section}{\normalfont\large\scshape\raggedright}{\thesection}{1em}{}
  \titleformat{name=\section,numberless}{\normalfont\large\scshape\raggedright}{}{0pt}{}
  \titleformat{\subsection}{\normalfont\normalsize\scshape\raggedright}{\thesubsection}{1em}{}
  \titleformat{name=\subsection,numberless}{\normalfont\normalsize\scshape\raggedright}{}{0pt}{}
  \titlespacing*{\section}{0pt}{2.0ex plus .4ex minus .2ex}{1.0ex}
  \titlespacing*{\subsection}{0pt}{1.7ex plus .3ex minus .2ex}{0.85ex}
  \newcommand{\papersection}[1]{\section{#1}}
  \newcommand{\papersubsection}[1]{\subsection{#1}}
  \newcommand{\paperacksection}{%
    \vspace{1.75em}%
    \par\noindent{\normalfont\normalsize\scshape Acknowledgments}\par\vspace{0.85em}%
  }
  \newcommand{\paperreferences}{%
    \vspace{1.25em}%
    \bibliography{assets/references}%
  }
  \justifying
  \AtBeginDocument{\captionsetup{justification=justified,singlelinecheck=false}}
\else
  \newcommand{\papersection}[1]{\section*{#1}}
  \newcommand{\papersubsection}[1]{\subsection*{#1}}
  \newcommand{\paperacksection}{\section*{Acknowledgments}}
  \newcommand{\paperreferences}{\bibliography{assets/references}}
  \raggedright
\fi
\newcommand{\paperfig}[2][\linewidth]{%
  \ifdraftfigures
    \par\medskip
    {\centering\includegraphics[width=#1,keepaspectratio]{#2}\par}%
    \par\medskip
  \fi
}
\newcommand{\suppfig}[2][0.88\linewidth]{%
  \ifdraftfigures
    \par\smallskip
    {\centering\includegraphics[width=#1,keepaspectratio]{#2}\par}%
    \par\smallskip
  \fi
}
\ifreportmode
  \newenvironment{paperfigure}{\begin{figure}[H]}{\end{figure}}
\else
  \newenvironment{paperfigure}{\begin{figure}[!ht]}{\end{figure}}
\fi

%\pagestyle{myheadings}
\pagestyle{fancy}
\fancyhf{}
%\setlength{\headheight}{27.023pt}
%\lhead{\includegraphics[width=2.0in]{PLOS-submission.eps}}
\rfoot{\thepage/\pageref{LastPage}}
\renewcommand{\headrulewidth}{0pt}
\renewcommand{\footrule}{\hrule height 2pt \vspace{2mm}}
\fancyheadoffset[L]{2.25in}
\fancyfootoffset[L]{2.25in}
\lfoot{\today}

%% END MACROS SECTION


\begin{document}
\ifreportmode
  \newgeometry{top=0.85in,left=1in,right=1in,footskip=0.75in}
  \fancyheadoffset[L]{0pt}
  \fancyfootoffset[L]{0pt}
  \fancyhead[C]{\small\textit{To be submitted to PLOS Computational Biology}}
\fi
\vspace*{0.08in}

\begin{center}
\ifreportmode
{\LARGE Rethinking Sleep Staging:\ }{\LARGE Unsupervised cross-laboratory mouse sleep staging with temporal variational autoencoders and substage discovery
\par\vspace{0.45em}}
\else
{\LARGE Rethinking Sleep Staging:\ }{\LARGE Unsupervised cross-laboratory mouse sleep staging with temporal variational autoencoders and substage discovery
\par\vspace{0.45em}}
\fi
\end{center}
\begin{center}
\begin{tabular*}{\textwidth}{@{\extracolsep{\fill}}c c c@{}}
\textbf{Oliver Rosb\ae{}k Elmgreen}\textsuperscript{1*} &
\textbf{Morten M\o{}rup}\textsuperscript{1} &
\textbf{Birgitte Rahbek Kornum}\textsuperscript{2} \\
\end{tabular*}
\end{center}
\begin{flushleft}
\textsuperscript{1} Department of Applied Mathematics and Computer Science, Technical University of Denmark
\\
\textsuperscript{2} Department of Neuroscience, University of Copenhagen, Denmark
\\
\smallskip
\textsuperscript{*} \texttt{s204070@dtu.dk}

\end{flushleft}

\begin{center}
\begin{minipage}{0.92\linewidth}
\begin{center}
\ifreportmode{\large\scshape Abstract}\else{\large\bfseries Abstract}\fi
\end{center}
\vspace{0.35em}
\justifying
\noindent
Cross-laboratory manual scoring of mouse sleep from electroencephalography (EEG) and electromyography (EMG) lacks consensus, and supervised models trained at one site often fail elsewhere without retraining or per-animal calibration.
We study twenty quality-screened wild-type mice from three European Mouse Sleep Staging Validation (MSSV) laboratories under leave-mice-out cross-validation with four folds.
Decoder-only conditional variational autoencoders (cVAE) assign states to held-out animals from spectral EEG and EMG alone, using only the encoder and a learned mixture prior at inference.
A hidden Markov model--Gaussian mixture model (HMM--GMM) prior that models bout structure reached mean prior normalized mutual information (NMI) of 0.575 on joint cross-laboratory holdout, versus 0.282 for a matched static mixture, a gain of 0.29 NMI averaged over folds.
Beyond three macro states, holdout macro NMI falls as mixture count $K$ increases while prior log-likelihood improves.
We therefore select $K$ for biological interpretability near six to nine substates rather than from macro NMI alone.
A population sweep on all twenty mice supports a $K{=}6$ map with six active substates and distinct Wake and non-rapid eye movement (NREM) components.
The method enables comparison of sleep substage structure across acquisition sites before expert labels are harmonised.
\end{minipage}
\end{center}

\section*{Author summary}
Sleep biologists score mouse EEG and EMG into Wake, NREM, and rapid eye movement (REM) sleep, but raters at different laboratories often disagree.
We trained unsupervised models on mice from Bern, Copenhagen, and Lyon, then tested them on held-out animals the models had never seen, without calibrating each new recording by hand.
Our best model reads spectral brain and muscle signals through a neural network and enforces realistic sleep bout structure with a temporal prior.
It outperformed simpler baselines on cross-laboratory holdout and recovered finer substates whose spectra and bout lengths matched patterns we discussed with sleep neuroscientists at the Kornum laboratory.
The approach is not yet accurate enough to replace supervised scoring, especially for REM, but it shows that substage structure can be studied across sites before expert labels are aligned.

\ifreportmode\else\newgeometry{top=0.85in,left=1in,right=1in,footskip=0.75in}\fi
\ifreportmode\else\linenumbers\fi

\papersection{Introduction}

Manual mouse sleep scores disagree across MSSV laboratories, and supervised models trained at one site fail on held-out animals elsewhere unless retrained or calibrated per mouse~\cite{rose2026consensus,miladinovic2019spindle,wang2026rest,barger2019accusleep,yaghouby2016segway}.
Expert staging nonetheless maps continuous physiology into Wake, NREM, and REM, the standard three-class benchmark.
Unsupervised work in mice has repeatedly resolved substages with distinct spectra, bout lengths, and transition roles that Wake/NREM/REM labels alone do not name~\cite{katsageorgiou2018mcRBM,stevner2019human,elmgreen2026thesis}.
Montage and laboratory protocols differ across studies~\cite{rose2025mssv}.
Claims that sleep structure extends beyond macro scoring must therefore be tested on new animals and new sites under population holdout without per-animal calibration at inference.

Unsupervised sequence models infer latent structure directly from EEG and EMG dynamics and align naturally with bout persistence and transitions~\cite{stevner2019human}.
Static clusterers and shallow hidden Markov models (HMMs) discover states from raw or hand-crafted features but rarely evaluate zero-shot transfer to held-out mice from other laboratories under a shared population model~\cite{sunagawa2013faster,yamada2024faster2,cusinato2024wucss}.
Katsageorgiou et al.\ resolved rodent substages beyond the macro taxonomy with distinct physiology and bout structure in mcRBM~\cite{katsageorgiou2018mcRBM}, yet without cross-laboratory leave-mice-out holdout.
Garc\'ia Ciudad et al.\ showed that longer temporal context improves supervised rodent staging, especially for REM~\cite{garciaciudad2024temporal}.
That setting remains fully supervised and does not test unsupervised cross-site transfer.
In a single-laboratory cohort with encoder+decoder conditioning and expert features, mixture-prior conditional Gaussian mixture variational autoencoder (cGMVAE) reached population prior NMI of about 0.70 under leave-one-subject-out prediction~\cite{elmgreen2026thesis}.
Cross-laboratory holdout on spectral convolutional neural network (CNN) inputs is substantially harder and is the focus here.

We extend decoder-only cVAE with HMM--GMM mixture priors~\cite{dilokthanakul2016gmvae,sohn2015cvae,johnson2016svae,ebbers2017hmmvae}.
Subject embeddings condition reconstruction during training only.
Held-out mice are staged from encoder means and the learned prior without subject identifiers.
We validate zero-shot cross-laboratory holdout at $K{=}3$, then test whether latents support finer structure than Wake, NREM, and REM through mixture-resolution sweeps and physiology-first interpretation~\cite{katsageorgiou2018mcRBM,elmgreen2026thesis}.
We compare a static Gaussian mixture model (GMM) prior with a warm conditional HMM--GMM variational autoencoder (cHMM--GMVAE) prior on locked spectral inputs across laboratories 2, 3, and 5 at Bern, Copenhagen, and Lyon.
Holdout fold~4 is the generalisation benchmark for macro metrics.
The full cohort supports substage physiology panels in the Appendix.
Dataset and protocol details are in Section~\ref{sec:data} and Appendix~A.

\papersection{Data}
\label{sec:data}

\methodpara{Dataset and cohort.}
MSSV on OpenNeuro ds006366~\cite{rose2025mssv} comprises invasive EEG and neck EMG from 92 mice across five European laboratories.
Expert Wake, NREM, REM, and Artifact labels are provided on a 4\,s grid for evaluation only.
We analyse laboratories 2 at Bern, 3 at Copenhagen, and 5 at Lyon.
These are healthy wild-type sites with different electrode layouts.
We exclude laboratories 1 and 4 because they are pathology-specific or single-electrode cohorts.
Quality screening retains mice with low EEG drift, defined as relative 0.5--2\,Hz power $\leq 0.2$ of 0.5--30\,Hz on each retained electrode.
EMG must separate Wake from NREM with a Wake/NREM root-mean-square ratio $\geq 2$.
The screened cohort comprises 20 mice with 6, 10, and 4 per site and ${\sim}453$\,h of recording under four-fold leave-mice-out cross-validation.
The subset balances cross-laboratory coverage with stable spectral inputs.
Training mice never appear in the validation fold for a given split.
No new animal procedures were performed.

\methodpara{Preprocessing and inputs.}
Montages are harmonised to a parietal--frontal EEG differential plus neck EMG.
Laboratory 2 uses EEG1 and EEG4.
Laboratories 3 and 5 use EEG1 and EEG2.
Per-site spectral filters and band limits were fixed from incohort ablations before holdout.
Recordings are segmented into 4\,s epochs at 128\,Hz.
Artifact epochs are removed.
Each epoch is a log-power spectral map $\mathbf{x}_t \in \mathbb{R}^{C \times F}$, batched into contiguous sequences for temporal models without hand-crafted expert features.

\papersection{Methods}
\label{sec:methods}

\papersubsection{Model and training}

Spectral EEG/EMG epochs are mapped to a low-dimensional latent space by a cVAE.
A subject embedding conditions reconstruction during training only.
At holdout, staging uses encoder means and a learned mixture prior without subject identifiers.
We compare two prior structures on matched inputs and cross-validation folds.
cGMVAE uses a static GMM with $T{=}1$.
cHMM--GMVAE uses a warm HMM--GMM sequence prior with $T{=}64$ epochs, about 4.3\,min, and $K{=}3$ for macro staging.
HMM--GMM variational autoencoder without subject conditioning (HMMGMVAE) provides a sequence-prior baseline in Table~\ref{table1}~\cite{ebbers2017hmmvae}.

\methodpara{Encoder-decoder.}
Training batches are contiguous windows $\mathbf{x}_{1:T} \in \mathbb{R}^{T \times C \times F}$ with $T=64$.
A convolutional encoder maps each epoch to a diagonal Gaussian posterior,
\begin{equation}
\label{eq:encoder}
q_{\boldsymbol{\phi}}(\mathbf{z}_t \mid \mathbf{x}_t) = \mathcal{N}\!\left(\mathbf{z}_t \mid \boldsymbol{\mu}_t, \operatorname{diag}(\boldsymbol{\sigma}_t^2)\right),
\qquad
(\boldsymbol{\mu}_t, \log \boldsymbol{\sigma}_t^2) = f_{\boldsymbol{\phi}}(\mathbf{x}_t),
\end{equation}
with $f_{\boldsymbol{\phi}}$ a CNN--multilayer perceptron (MLP) encoder and $\mathbf{z}_t \in \mathbb{R}^{d}$, $d{=}6$.
Latent samples use the reparameterisation $\mathbf{z}_t = \boldsymbol{\mu}_t + \boldsymbol{\sigma}_t \odot \boldsymbol{\epsilon}_t$ with $\boldsymbol{\epsilon}_t \sim \mathcal{N}(\mathbf{0}, \mathbf{I})$.
The decoder reconstructs each epoch from $\mathbf{z}_t$ and, during training only, a subject embedding $\mathbf{e}_i \in \mathbb{R}^{8}$:
\begin{equation}
\label{eq:decoder}
\hat{\mathbf{x}}_t = g_{\boldsymbol{\theta}}\!\left(\mathbf{z}_t, \mathbf{e}_i\right).
\end{equation}
Under decoder-only conditioning, $f_{\boldsymbol{\phi}}$ does not receive $\mathbf{e}_i$.
CNN and MLP widths and the full architecture diagram are in \apdx{app:methods:hyper} and \nameref{fig_arch}.

\methodpara{Training objective.}
For each window $\mathbf{x}_{1:T}$, training minimises
\begin{equation}
\label{eq:loss}
\mathcal{L} = \mathcal{L}_{\mathrm{recon}} + \beta\,\mathcal{L}_{\mathrm{KL}},
\end{equation}
where
\begin{equation}
\label{eq:recon}
\mathcal{L}_{\mathrm{recon}} = \frac{1}{T}\sum_{t=1}^{T}\left\|\mathbf{x}_t - \hat{\mathbf{x}}_t\right\|_2^2
\end{equation}
is mean squared reconstruction error over $T$ contiguous 4\,s epochs.
The Kullback--Leibler (KL) term is the sequence-averaged log-density ratio between the encoder posterior and the mixture prior,
\begin{equation}
\label{eq:kl_hmm}
\mathcal{L}_{\mathrm{KL}} = \frac{1}{T}\sum_{t=1}^{T}\left[\log q_{\boldsymbol{\phi}}(\mathbf{z}_t \mid \mathbf{x}_t) - \log p(\mathbf{z}_{1:T})\right].
\end{equation}
For cHMM--GMVAE, $\log p(\mathbf{z}_{1:T})$ is the HMM--GMM forward marginal with emission densities $b_k(\mathbf{z}_t)$ and transition matrix $A$.
For cGMVAE, each epoch is scored independently with a static GMM marginal at $T{=}1$.
The weight $\beta$ is zero for the first ten epochs, then ramps linearly from 0.01 to 1.0.
cHMM--GMVAE warm-starts the mixture prior from encoder means~\cite{ebbers2017hmmvae,johnson2016svae}.
Phase timings and prior equations are in \apdx{S1_Appendix} and \apdx{app:methods:warmstart}.

\methodpara{Prior-based staging at holdout.}
Held-out mice are staged from encoder means $\boldsymbol{\mu}_t$ and the trained prior without subject embeddings or decoder reconstruction.
Static cGMVAE assigns $\hat{s}_t = \arg\max_k \left[\log \pi_k + \log b_k(\boldsymbol{\mu}_t)\right]$.
cHMM--GMVAE applies Viterbi decoding,
\begin{equation}
\label{eq:viterbi}
\hat{\mathbf{s}}_{1:T} = \arg\max_{s_{1:T}} \left[\log \pi_{s_1} + \sum_{t=2}^{T}\log A_{s_{t-1}, s_t} + \sum_{t=1}^{T}\log b_{s_t}(\boldsymbol{\mu}_t)\right],
\end{equation}
with $B_{tk} = \log b_k(\boldsymbol{\mu}_t)$.

\papersubsection{Evaluation protocol}

The twenty quality-screened mice were partitioned into four leave-mice-out folds (\apdx{app:data:splits}).
Joint scope pools all non-held-out mice from laboratories 2, 3, and 5 into one population model and tests on the held-out fold across sites.
Within-lab scope restricts both training and validation to a single laboratory with the same fold assignments.
Held-out mice are staged from encoder means and the trained mixture prior only, as in Eqs.~\eqref{eq:viterbi} and the holdout paragraph above.
Per-laboratory spectral locks remain fixed.
Agreement with expert Wake, NREM, and REM is prior NMI,
\begin{equation}
\label{eq:nmi}
\mathrm{NMI}(Y,\hat{Y}) = \frac{2\,I(Y;\hat{Y})}{H(Y)+H(\hat{Y})},
\end{equation}
where $I$ is mutual information and $H$ Shannon entropy.
Each fold trains three seeds; holdout prior NMI uses the best unsupervised checkpoint per seed and reports the maximum across seeds (Appendix~B).
Macro-aligned accuracy, Cohen's $\kappa$, decoder-only ablation, and substage sweep setup are in Appendix~B and~C.

\methodpara{Data and code availability.}
MSSV is available on OpenNeuro as ds006366.
Processed features, analysis code, and full experimental settings are provided as supplementary material in \nameref{S1_File}.

\papersection{Results and discussion}

\papersubsection{Cross-laboratory holdout at macro resolution}

We evaluated decoder-only cGMVAE and cHMM--GMVAE under joint leave-mice-out training on twenty quality-screened mice from Bern, Copenhagen, and Lyon (Table~\ref{table1}).
Supervised pipelines on MSSV reach much higher accuracy when retrained per site or given long temporal context~\cite{miladinovic2019spindle,wang2026rest,garciaciudad2024temporal}.
Under the same unsupervised constraint, cHMM--GMVAE reached mean prior NMI 0.575 versus 0.282 for matched cGMVAE, a gain of 0.29 averaged over four folds.

\begin{table}[!ht]
\centering
\caption{Joint holdout prior NMI by fold (best-of-three seeds, locked recipes, $K{=}3$).
The raw time-domain hidden Markov model (HMM) baseline folds 1--3 report best logged validation prior NMI from incomplete joint runs with $\leq 0.003$.
Fold~4 uses the full holdout run.}
% Auto-generated by scripts/paper/build_holdout_scope_table.py
\begin{tabular}{|l|r|r|r|r|}
\hline
\textbf{Fold} & \textbf{HMM (raw)} & \textbf{cGMVAE} & \textbf{HMMGMVAE} & \textbf{cHMM--GMVAE} \\
\thickhline
1 & 0.003 & 0.267 & 0.409 & 0.579 \\
2 & 0.001 & 0.256 & 0.236 & 0.396 \\
3 & 0.003 & 0.307 & 0.497 & 0.675 \\
4 & 0.208 & 0.299 & 0.404 & 0.649 \\
\hline
\textbf{Mean} & 0.054 & 0.282 & 0.387 & 0.575 \\
\hline
\end{tabular}

\label{table1}
\end{table}

In Table~\ref{table1}, the temporal prior improves every fold.
The largest margins appear on folds 3 and 4 at +0.36 and +0.41 NMI.
The smallest gain is on fold 2 at +0.10, where Bern montages mix with held-out mice from all sites.
Fold~4 confusion (\nameref{S3_Fig}) shows strong Wake and NREM alignment but REM recall near 3\%.
Thus unsupervised models recover coarse Wake/NREM/REM structure on held-out mice across laboratories, well above chance but below supervised ceilings.
Three macro classes remain an incomplete description of the signal.

\papersubsection{Laboratory-specific transfer}

Table~\ref{table2} compares joint population training with within-laboratory training on the \emph{same} held-out mice from each site.
For joint training, NMI and macro-aligned accuracy are computed on laboratory-specific holdout subsets only, not pooled across sites.

\begin{table}[!ht]
\centering
\caption{Fair per-laboratory holdout comparison (mean over four folds, $K{=}3$).
Joint training uses all laboratories but reports metrics on held-out mice from the listed site only.
Within training trains and evaluates inside the same site.
All rows use locked preprocessing per laboratory and matched CNN--variational autoencoder (VAE) capacity with $d{=}6$, MLP widths 512--256--128, 200 epochs, batch 128, and three seeds.
cGMVAE, HMMGMVAE, and cHMM--GMVAE share learning rate $1.3 \times 10^{-3}$ and differ in embedding use, sequence length, and prior type.
HMM (raw) uses the time-domain pipeline without a shared deep front-end.
VAE models report best-of-three prior NMI and matching macro accuracy, averaged across folds.
HMM (raw) has $\mathit{runs}{=}1$.}
% Auto-generated by scripts/paper/build_holdout_scope_table.py
\begin{tabular}{|l|l|r|r|r|r|}
\hline
\textbf{Lab} & \textbf{Model} & \textbf{Joint NMI} & \textbf{Joint Acc.} & \textbf{Within NMI} & \textbf{Within Acc.} \\
\thickhline
Lab 2 & HMM (raw) & 0.179 & 0.665 & 0.000 & 0.495 \\
 & cGMVAE & 0.224 & 0.648 & 0.549 & 0.855 \\
 & HMMGMVAE & 0.039 & 0.543 & 0.586 & \textbf{0.865} \\
 & cHMM--GMVAE & \textbf{0.570} & \textbf{0.869} & \textbf{0.596} & 0.864 \\
\hline
Lab 3 & HMM (raw) & 0.237 & 0.731 & 0.004 & 0.475 \\
 & cGMVAE & 0.323 & 0.727 & 0.550 & 0.850 \\
 & HMMGMVAE & 0.586 & 0.850 & 0.637 & 0.897 \\
 & cHMM--GMVAE & \textbf{0.705} & \textbf{0.923} & \textbf{0.738} & \textbf{0.939} \\
\hline
Lab 5 & HMM (raw) & 0.240 & 0.724 & 0.002 & 0.525 \\
 & cGMVAE & 0.245 & 0.700 & 0.490 & 0.856 \\
 & HMMGMVAE & 0.483 & 0.821 & \textbf{0.584} & \textbf{0.882} \\
 & cHMM--GMVAE & \textbf{0.556} & \textbf{0.877} & 0.565 & 0.881 \\
\hline
\end{tabular}

\label{table2}
\end{table}

Within-lab training exceeds joint training at every site.
On held-out Bern mice, cGMVAE prior NMI averages 0.22 under joint training versus 0.55 within-laboratory on the same animals.
The pattern is consistent with residual covariate shift when pooling heterogeneous spectral profiles across sites.
Decoder-only subject embeddings absorb amplitude nuisance during reconstruction but do not fully remove equipment and scoring context.

\papersubsection{Substage resolution and physiology}

Expert Wake, NREM, and REM are a post-hoc reference for NMI, not a claim that sleep contains only three dynamical regimes.
On holdout fold~4, macro NMI peaks near $K{=}3$--$4$ and falls at higher $K$, while prior log-likelihood continues to rise (Fig~\ref{fig2}A).
The same trade-off appears on the full cohort (Fig~\ref{fig2}B).
Fold~4 alone is under-powered for stable substage naming at $K \geq 6$.

\begin{paperfigure}
\noindent
\begin{minipage}[t]{0.48\linewidth}
\paperfig[\linewidth]{figures/main/S3_k_sweep.pdf}
\centerline{\small A Holdout fold~4}
\end{minipage}\hfill
\begin{minipage}[t]{0.48\linewidth}
\paperfig[\linewidth]{figures/supplementary/S15_population_k_sweep.pdf}
\centerline{\small B All 20 mice (population)}
\end{minipage}
\caption{Substage resolution sweeps.
Mean validation prior NMI (blue) and mean normalised prior log-density (orange) versus mixture count $K$. Shaded bands: mean $\pm$ standard deviation (SD) across seeds.
A, Joint leave-mice-out holdout on fold~4 (six mice).
B, Population sweep on the full cohort (same locked recipe).
Macro NMI and prior density diverge once $K$ exceeds the expert three-class grid.}
\label{fig2}
\end{paperfigure}

We therefore interpret substage maps from a population sweep on all twenty mice under the locked recipe.
Expert labels enter only for post-hoc NMI and physiology review.
Fig~\ref{fig3} shows clearer latent separation at $K{=}6$ and $K{=}9$ than at $K{=}12$, with brief boundary excursions visible on the hypnogram.
State distinctness in \nameref{S12_Table} provides a label-free backup when macro NMI falls.
At population $K{=}6$, all six components are active with $p_{\mathrm{perm}} = 0.02$, Cram\'{e}r's $V = 0.61$, and energy distance (ED) $= 2.49$ despite macro NMI $= 0.46$.
Latent substates therefore remain separated even when macro agreement falls.

\begin{paperfigure}
\paperfig[\linewidth]{figures/main/Fig_population_pca_2x2.png}
\caption{Latent separation and temporal structure (population sweep, best seed).
A--B, principal component analysis (PCA) of encoder means: expert Wake/NREM/REM versus predicted substates at $K{=}6$ and $K{=}9$.
C, Finer latent substructure at $K{=}12$.
D, Example 30\,min hypnogram at $K{=}12$ with macro-group background bands.}
\label{fig3}
\end{paperfigure}

Fig~\ref{fig4} summarises substage physiology at $K{=}6$.
\textbf{S2} is core quiet Wake with long bouts and about 98\% expert Wake.
\textbf{S5} and \textbf{S6} are dominant NREM substates with separated depth in spectra.
\textbf{S3} consolidates REM at about 66\% expert REM and does not split further even at $K{=}12$~\cite{katsageorgiou2018mcRBM,elmgreen2026thesis}.
\textbf{S1} and \textbf{S4} are short-bout boundary substates between Wake and NREM.
Transition structure at $K{=}6$ is in \nameref{fig5}.
Together, these panels support a Wake/NREM-plus map near six substates rather than choosing $K$ from the macro-NMI peak alone.

\begin{paperfigure}
\paperfig[0.95\linewidth]{figures/main/Fig4_K6_frequency.png}
\caption{Substage physiology at $K{=}6$ (population sweep, best seed).
Per-substage EEG/EMG power spectra, $\theta/\delta$, bout length, expert macro composition, and subject mix.
Wake and NREM split into substates with distinct spectra; REM remains a single consolidated class.}
\label{fig4}
\end{paperfigure}

\papersubsection{Limitations}
\label{sec:limitations}

Agreement is quantified by NMI against expert macro labels, which inherit cross-laboratory disagreement~\cite{rose2026consensus}.
REM recall on fold~4 remains near 3\% despite temporal context.
Cross-laboratory holdout is harder than within-laboratory training on the same mice (Table~\ref{table2}).
The holdout substage sweep uses fold~4 only.
Population $K{=}6$ panels train on all twenty mice and support physiology, not leave-mice-out generalisation claims.
Reported NMI uses best-of-three seeds per fold, which is optimistic relative to seed means.
The method is not yet accurate enough to replace supervised staging, especially for REM.

\papersection{Conclusion}

Cross-laboratory mouse sleep can be staged without labels or per-mouse calibration, but only coarsely today.
Temporal mixture priors improve Wake/NREM alignment on held-out mice relative to a static mixture prior, with mean prior NMI 0.575 versus 0.282 (+0.29 over folds).
On the full cohort, finer structure appears at $K{=}6$, with six active substates separating Wake and NREM components.
State distinctness remains high at $K{=}6$ even as macro NMI falls, which supports reading low NMI as label mismatch rather than collapsed mixtures.
The approach can support comparison across acquisition sites before expert labels are harmonised.
It does not yet replace supervised scoring.

\paperacksection
This report extends an MSc thesis co-authored with Andreas Holmer Bigom. Thank you to Javier Garc\'ia Ciudad and Alberte Wollesen J\"urgensen-Breum for machine-learning and sleep-biology guidance.

\nolinenumbers

\paperreferences

\clearpage
\appendix
\justifying

\section{Data and cohort}
\label{app:data}
% Appendix A — Data (first appendix section).
\subsection{MSSV cohort and quality screening}
\label{app:data:cohort}

MSSV (OpenNeuro ds006366)~\cite{rose2025mssv} comprises invasive EEG and neck EMG from 92 mice across five European laboratories, with expert Wake, NREM, REM, and Artifact labels on a 4\,s grid for evaluation only.
This report analyses laboratories 2 (Bern), 3 (Copenhagen), and 5 (Lyon): healthy wild-type sites with different electrode layouts, excluding laboratories 1 and 4 (pathology-specific or single-electrode cohorts).

Quality screening retains mice with low EEG drift and EMG that separates Wake from NREM.
Drift is quantified as relative 0.5--2\,Hz power $\leq 0.2$ of 0.5--30\,Hz on each retained electrode.
Wake/NREM separability uses a root-mean-square ratio $\geq 2$ on neck EMG.
The screened cohort comprises 20 mice (6 Bern, 10 Copenhagen, 4 Lyon) with 42 recording runs and ${\sim}453$\,h of EEG/EMG (${\sim}4.1 \times 10^{5}$ artefact-free 4\,s epochs after screening).
Training mice never appear in the validation fold for a given split (\apdx{app:data:splits}).
No new animal procedures were performed.

\apptabdesc{Quality-screened cohort inventory. Per-mouse recording runs, hours, and 4\,s epoch counts (MSSV annotation grid).}
\label{app:tab:cohort}
\apptabular{\small% Auto-generated by scripts/paper/build_cv4fold_split_tables.py
\begin{tabular}{|l|l|r|r|r|}
\hline
\textbf{Mouse} & \textbf{Site} & \textbf{Runs} & \textbf{Hours} & \textbf{Epochs (4\,s)} \\
\thickhline
sub-071 & Bern (2) & 2 & 12 & 10{,}800 \\
sub-072 & Bern (2) & 2 & 12 & 10{,}800 \\
sub-076 & Bern (2) & 2 & 12 & 10{,}800 \\
sub-077 & Bern (2) & 2 & 13 & 11{,}700 \\
sub-080 & Bern (2) & 2 & 12 & 10{,}800 \\
sub-081 & Bern (2) & 2 & 12 & 10{,}800 \\
sub-038 & Copenhagen (3) & 2 & 48 & 43{,}201 \\
sub-039 & Copenhagen (3) & 3 & 12 & 10{,}836 \\
sub-041 & Copenhagen (3) & 3 & 12 & 10{,}803 \\
sub-043 & Copenhagen (3) & 1 & 24 & 21{,}600 \\
sub-048 & Copenhagen (3) & 3 & 12 & 10{,}804 \\
sub-054 & Copenhagen (3) & 2 & 48 & 43{,}194 \\
sub-056 & Copenhagen (3) & 1 & 23.6 & 21{,}196 \\
sub-059 & Copenhagen (3) & 1 & 24 & 21{,}600 \\
sub-060 & Copenhagen (3) & 1 & 8 & 7{,}212 \\
sub-069 & Copenhagen (3) & 1 & 24 & 21{,}600 \\
sub-087 & Lyon (5) & 6 & 72 & 64{,}788 \\
sub-088 & Lyon (5) & 2 & 24 & 21{,}595 \\
sub-089 & Lyon (5) & 2 & 24 & 21{,}596 \\
sub-092 & Lyon (5) & 2 & 24 & 21{,}595 \\
\hline
\textbf{Total (20 mice)} &  & 42 & 452.6 & 407{,}320 \\
\hline
\end{tabular}
}

\subsection{Preprocessing and spectral inputs}
\label{app:data:preprocessing}

Montages are harmonised to a parietal-frontal EEG differential plus neck EMG.
Bern (laboratory~2) provides four EEG contacts (EEG1/2 parietal, EEG3/4 frontal); we use EEG1 and EEG4 to match the single parietal--frontal pairing at laboratories~3 and~5 (EEG1 parietal, EEG2 frontal).
Recordings are segmented into 4\,s epochs at 128\,Hz. Artifact epochs are removed.
Each epoch is a log-power spectral map $\mathbf{x}_t \in \mathbb{R}^{C \times F}$, batched into contiguous sequences for temporal models without hand-crafted expert features.

Per-laboratory channel sets were fixed from incohort ablation winners (full sweep outcomes in \apdx{app:incohort_ceiling}).
Spectral passbands are summarised in \apdx{app:data:bands}.

\subsection{Spectral bands and frequency content}
\label{app:data:bands}

Each 4\,s epoch comprises 512 samples at 128\,Hz sampling rate.
The CNN input is a log-power spectrum from an unwindowed real FFT (frequency resolution 0.25\,Hz; Nyquist 64\,Hz).
Band limiting is applied in the frequency domain on the complex spectrum before log-power: a null lower cutoff retains all bins from DC upward, and an upper cutoff zeroes bins above the configured limit.
The model therefore receives the full spectral shape inside each passband, not a small vector of hand-defined sleep bands.

\paragraph*{Locked per-laboratory passbands (holdout protocol).}
\apptabular{\small
\begin{tabular}{|l|l|l|l|}
\hline
\textbf{Site} & \textbf{EEG channels} & \textbf{EEG passband} & \textbf{EMG passband} \\
\thickhline
Bern (lab~2) & EEG1, EEG4 & 0--30\,Hz & 3--100\,Hz ($\leq 64$\,Hz Nyquist) \\
Copenhagen (lab~3) & EEG1, EEG2 & 0--20\,Hz & 5--60\,Hz \\
Lyon (lab~5) & EEG1, EEG2 & 0--20\,Hz & 5--60\,Hz \\
\hline
\end{tabular}}

\paragraph*{Diagnostic and physiology bands.}
Quality screening compares relative 0.5--2\,Hz EEG power to 0.5--30\,Hz total (\apdx{app:data:cohort}).
For substage physiology panels (Fig~\ref{fig3} and Fig~\ref{fig4}) and summary tables, log-power spectra are aggregated into interpretable bands on the same frequency grid:
delta 0.5--4\,Hz, theta 4--8\,Hz, and broadband EMG power summed over the retained EMG passband.
Reported $\theta/\delta$ ratios use these bands.

Bern uses no post-normalisation on log-power spectra.
Copenhagen and Lyon apply run-wise z-scoring after band limiting: for each recording run, each channel's log-power spectrum is scaled by the mean and standard deviation computed across all 4\,s epochs in that run (post-normalisation).
These limits were chosen from incohort ablations (\apdx{app:incohort_ceiling}) and held fixed for all cross-laboratory holdout runs.

\subsection{Cross-validation protocol}
\label{app:data:splits}

Folds were designed so each site contributes holdout mice every fold while balancing validation load: fold~1 is epoch-heaviest (${\sim}32\%$ of cohort epochs in validation, driven by long Copenhagen and Lyon recordings) whereas folds~2--4 hold out 21--25\% of epochs each.
We use four folds because Lyon contributes only four quality-screened mice, which is the maximum fold count for strict mouse-level holdout.
sub-069 is assigned to fold~2 rather than fold~4 so Copenhagen does not contribute four simultaneous holdout mice on the substage-analysis fold.
Each VAE configuration trains three random seeds per fold. Reported metrics use the best prior NMI among seeds.

\apptabdesc{Four-fold leave-mice-out assignments. Top: held-out mice per site and fold.
Bottom: validation versus training pool sizes (joint scope trains on all non-held-out mice across sites).}
\label{app:tab:splits}
\apptabular{% Auto-generated by scripts/paper/build_cv4fold_split_tables.py
\begin{tabular}{|c|l|l|r|r|r|}
\hline
\textbf{Fold} & \textbf{Site} & \textbf{Held-out mice} & \textbf{Runs} & \textbf{Hours} & \textbf{Epochs} \\
\thickhline
1 & Bern (2) & sub-071 & 2 & 12 & 10{,}800 \\
 & Copenhagen (3) & sub-038, sub-039 & 5 & 60 & 54{,}037 \\
 & Lyon (5) & sub-087 & 6 & 72 & 64{,}788 \\
\hline
2 & Bern (2) & sub-072, sub-076 & 4 & 24 & 21{,}600 \\
 & Copenhagen (3) & sub-041, sub-048, sub-069 & 7 & 48 & 43{,}207 \\
 & Lyon (5) & sub-088 & 2 & 24 & 21{,}595 \\
\hline
3 & Bern (2) & sub-077 & 2 & 13 & 11{,}700 \\
 & Copenhagen (3) & sub-043, sub-054 & 3 & 72 & 64{,}794 \\
 & Lyon (5) & sub-089 & 2 & 24 & 21{,}596 \\
\hline
4 & Bern (2) & sub-080, sub-081 & 4 & 24 & 21{,}600 \\
 & Copenhagen (3) & sub-056, sub-059, sub-060 & 3 & 55.6 & 50{,}008 \\
 & Lyon (5) & sub-092 & 2 & 24 & 21{,}595 \\
\hline
\end{tabular}
\par\medskip
\begin{tabular}{|c|r|r|r|r|r|r|r|}
\hline
\textbf{Fold} & \textbf{Val mice} & \textbf{Val runs} & \textbf{Val hours} & \textbf{Val epochs} & \textbf{Train mice} & \textbf{Train epochs} & \textbf{Val (\%)} \\
\thickhline
1 & 4 & 13 & 144 & 129{,}625 & 16 & 277{,}695 & 31.8 \\
2 & 6 & 13 & 96 & 86{,}402 & 14 & 320{,}918 & 21.2 \\
3 & 4 & 7 & 109 & 98{,}090 & 16 & 309{,}230 & 24.1 \\
4 & 6 & 9 & 103.5 & 93{,}203 & 14 & 314{,}117 & 22.9 \\
\hline
\end{tabular}
}


\section{Training, checkpoints, and metrics}
\label{S1_Appendix_methods}
% Technical methods detail moved from main text for readability (Appendix B).

\subsection{Locked hyperparameters}
\label{app:methods:hyper}
Latent dimension $d{=}6$ (matched to thesis-scale mixture models), embedding dimension 8, learning rate $1.3 \times 10^{-3}$, batch size 128, 200 epochs, three random seeds per fold.
Sequence length $T{=}64$ epochs (${\sim}4.3$\,min); sticky self-transition weight $\kappa{=}0.92$.
CNN channel widths $[32,64,128,128]$, kernel sizes $[7,5,5,3]$, strides 2; MLP widths $[512,256,128]$ (encoder) and $[128,256,512]$ (decoder).

\subsection{Incohort ablation rationale}
Per-laboratory spectral bandpass, montage, post-normalisation, and MLP width were fixed from incohort sweeps before any holdout run (\apdx{app:data:preprocessing}, \apdx{app:incohort_ceiling}).
These runs pool all mice within a site and estimate a performance \emph{ceiling}; they do not use leave-mice-out folds.

\subsection{Sensitivity checks}
On joint holdout fold~4, sequence length $T{=}64$ with warm sticky training reached prior NMI ${\sim}0.65$ (best seed), compared with ${\sim}0.17$ for $T{=}32$ without the warm schedule.
On laboratory~3 incohort, warm-started HMM--GMM reached best-of-three prior NMI 0.734 versus 0.694 for a cold HMM--GMM start.
Locked settings prioritise reproducibility across holdout runs.

\subsection{Warm prior schedule (cHMM--GMVAE)}
\label{app:methods:warmstart}
Epochs 1--17: standard Gaussian KL on the encoder.
Epoch 18: $K$-means on encoder means initialises GMM emissions.
Epochs 18--36: independent mixture KL.
From epoch 37: sticky HMM--GMM transitions ramp in over 18 epochs.
$\beta$ is 0 for the first ten epochs, then ramps linearly from 0.01 to 1.0.

\subsection{Staging on encoder means}
At holdout, emissions use $\boldsymbol{\mu}_t$ rather than sampled $\mathbf{z}_t$.
The posterior mean can bias emission scores slightly relative to stochastic latents but reduces Monte Carlo variance in Viterbi decoding and improves reproducibility across seeds once the encoder is trained.

\subsection{Unsupervised checkpoint score}
Expert labels are not used to select weights during training.
Let $\tilde{\omega}(\varepsilon)$ denote validation log-likelihood (sequence-normalised negative ELBO) and $H_{\mathrm{norm}}(\hat{Y})$ the entropy of the predicted state histogram normalised by $\log K$ ($0$ = collapse to one state, $1$ = uniform occupancy).
Following the thesis composite criterion~\cite{elmgreen2026thesis},
\begin{equation}
\label{eq:checkpoint}
S(\varepsilon) = \beta_{\mathrm{ckpt}}\,\tilde{\omega}(\varepsilon) + (1-\beta_{\mathrm{ckpt}})\,H_{\mathrm{norm}}(\hat{Y}),
\qquad \beta_{\mathrm{ckpt}} = 0.9.
\end{equation}
The epoch maximising $S(\varepsilon)$ after a 5\% warmup is retained per seed.
Prior NMI against expert macros is computed only \emph{after} checkpoint selection for reporting.

\subsection{Macro-aligned accuracy and label switching}
Macro accuracy maps each predicted cluster to its dominant expert macro label (Hungarian alignment when $K{=}3$), then reports the fraction of epochs matching Wake/NREM/REM.
Accuracy can exceed chance and diverge from NMI when $K>3$ substates collapse onto the same macro label, when too few clusters are active, or when clusters are macro-pure but permuted.
Mixture components are identifiable only up to permutation; NMI is invariant to permutations, but cluster indices differ across seeds unless post-hoc alignment is applied for tables and figures.

\subsection{Raw time-domain HMM baseline}
\label{app:methods:hmm_raw}
Classical HMMs on raw MSSV EEG/EMG voltages provide a time-domain reference without a shared spectral CNN front-end.
On the locked cross-lab protocol, joint HMM (raw) prior NMI averaged 0.054 across four folds (Table~\ref{table1}), motivating the spectral deep-latent models used throughout this report.

\subsection{Seed aggregation}
Each fold trains three seeds; holdout prior NMI is evaluated at each seed's $S(\varepsilon)$ checkpoint; the maximum across seeds is the headline fold score (Tables~\ref{table1}--\ref{table2}).
All four folds favour cHMM--GMVAE over cGMVAE under the locked recipe.

\subsection{Prior log-likelihood in K-sweeps}
K-sweep panels (Fig~\ref{fig2}) report mean normalised prior log-density alongside prior NMI.
For cGMVAE ($T{=}1$), density is per 4\,s epoch; for HMMGMVAE and cHMM--GMVAE ($T{=}64$), density is the HMM forward marginal per sequence window, normalised by $T \cdot L$ for comparability on the dual-axis plots.


\section{Prior models and protocol ablations}
\label{app:ablations}

\subsection{HMM--GMM prior (equations)}
\label{S1_Appendix}

Each discrete latent state $s_t \in \{1,\ldots,K\}$ emits a diagonal Gaussian in $\mathbb{R}^{d}$.
The per-step emission log-density is
\begin{equation}
\label{eq:emission}
\log b_k(\mathbf{z}_t) = \log \mathcal{N}\!\left(\mathbf{z}_t \mid \mathbf{m}_k, \operatorname{diag}(\mathbf{v}_k)\right).
\end{equation}
Sticky transitions use self-transition weight $\kappa{=}0.92$.
The marginal log-likelihood of a latent window is
\begin{equation}
\label{eq:forward}
\log p(\mathbf{z}_{1:T}) = \log \sum_{s_{1:T}} \pi_{s_1} \prod_{t=2}^{T} A_{s_{t-1}, s_t} \prod_{t=1}^{T} b_{s_t}(\mathbf{z}_t),
\end{equation}
which enters the training KL term $\mathcal{L}_{\mathrm{KL}}$ in Eq.~\eqref{eq:kl_hmm}.
Warm-start phases and CNN widths are in \apdx{app:methods:warmstart} and \apdx{app:methods:hyper}.
Controlled synthetic datasets in the MSc thesis validated MAR-HMM optimisation before MSSV experiments (thesis Ch.~3); that validation is not repeated here.

\subsection{Decoder-only vs encoder+decoder conditioning}
\label{S3_Appendix}

The MSc thesis default fed subject embeddings to \emph{both} encoder and decoder on laboratory~3 only~\cite{elmgreen2026thesis}.
Before holdout, we compared decoder-only and encoder+decoder paths on locked incohort recipes at all three sites ($K{=}3$, three seeds):
% Decoder-only vs encoder+decoder conditioning (incohort ablation).
\begin{center}
\begin{tabular}{|l|r|r|l|}
\hline
\textbf{Site} & \textbf{Decoder-only} & \textbf{Encoder+decoder} & \textbf{Best-of-three prior NMI} \\
\thickhline
Bern (lab~2) & 0.593 & 0.578 & Decoder-only \\
Copenhagen (lab~3) & 0.737 & 0.621 & Decoder-only \\
Lyon (lab~5) & 0.534 & 0.430 & Decoder-only \\
\hline
\end{tabular}
\end{center}

\medskip
\noindent
All three sites favoured decoder-only conditioning on incohort prior NMI.
The thesis headline (${\sim}0.70$ population NMI) used encoder+decoder on laboratory~3 with expert features and leave-one-subject-out evaluation~\cite{elmgreen2026thesis}, which is not directly comparable to these three-lab spectral incohort runs.

Holdout fold-level encoder+decoder comparisons were not run under the locked paper protocol; the main line uses decoder-only throughout.

% Within-laboratory ablation ceiling (locks preprocessing before holdout).
\subsection{Incohort performance ceiling}
\label{app:incohort_ceiling}

Before cross-laboratory holdout, we ran extensive within-laboratory sweeps over preprocessing, CNN--VAE architecture, prior tier (static GMM vs.\ HMM--GMM vs.\ warm sticky HMM--GMM), and sequence length $T$ (Table~\ref{app:tab:incohort}).
All mice in a laboratory were pooled for training and validation tuning (no leave-mice-out).
Reported values are best-of-three prior NMI at $K{=}3$ with decoder-only conditioning.
These runs fixed the per-site recipes in \apdx{app:data:preprocessing} and estimate an upper performance bound when site-specific tuning is allowed.
They are \emph{not} the locked holdout protocol in Table~\ref{table2}.
Macro accuracy was not logged systematically in the incohort sweeps.

Bern (laboratory~2): cGMVAE reached 0.593 with EEG1+EEG4, no post-normalisation, and EEG 0--30\,Hz.
Temporal priors did not beat the static mixture at $T{=}1$ (best cHMM--GMVAE 0.571 with warm prior and 50\,Hz notch).
Copenhagen (laboratory~3): cGMVAE peaked at 0.737; warm cHMM--GMVAE at $T{=}64$ matched within 0.003 (0.734).
Lyon (laboratory~5): cHMM--GMVAE with EMG-wide notch preprocessing at $T{=}32$ reached 0.640 versus 0.534 for cGMVAE.

\apptabdesc{Incohort ceiling prior NMI. Best-of-three prior NMI and winning configuration per laboratory and model.}
\label{app:tab:incohort}
\apptabular{% Incohort ceiling reference (best-of-3 prior NMI from per-lab ablation sweeps).
% Protocol: all mice in lab, no cross-fold holdout; per-lab preprocessing/arch locks.
% See docs/cv4fold/ablations/ablation_findings_20260608.md
\begin{tabular}{|l|l|r|l|}
\hline
\textbf{Lab} & \textbf{Model} & \textbf{Best NMI} & \textbf{Winning lock (summary)} \\
\thickhline
2 & cGMVAE & 0.593 & EEG1+EEG4, no post-normalisation, EEG 0--30\,Hz, EMG 3--100\,Hz \\
2 & cHMM--GMVAE & 0.571 & warm prior + 50\,Hz notch, below cGMVAE at $T{=}1$ \\
\hline
3 & cGMVAE & 0.737 & post-normalisation, EEG 0--20\,Hz, MLP widths 512--256--128 \\
3 & cHMM--GMVAE & 0.734 & warm HMM--GMM, $T{=}64$, same preprocessing and architecture locks \\
\hline
5 & cGMVAE & 0.534 & post-normalisation, EEG 0--20\,Hz, MLP widths 512--256--128 \\
5 & cHMM--GMVAE & 0.640 & EMG wide + notch, $T{=}32$ (1/3 seeds healthy) \\
\hline
\end{tabular}
}

\noindent{\small\textbf{Ceiling reference only.}
Holdout within-lab NMI and accuracy under the same locked recipes appear in Table~\ref{table2}.
Joint holdout remains below these ceilings because leave-mice-out validation is strictly harder than pooling all laboratory mice during tuning.
The ceilings also show that cGMVAE can reach much higher prior NMI under site-tuned incohort settings (for example 0.737 at Copenhagen) than under the unified holdout recipe (joint mean 0.282).
That gap is expected because holdout locks one preprocessing and training recipe across cGMVAE, HMMGMVAE, and cHMM--GMVAE for fair prior comparison rather than maximising the static mixture alone.
A useful follow-up is to select cGMVAE-specific hyperparameters from incohort sweeps at $T{=}1$ and re-run joint holdout, to see how much of the incohort ceiling transfers under leave-mice-out and whether cHMM--GMVAE retains its margin.}\par\vspace{0.5em}


\section{Supplementary figures}
\label{app:figures}

\subsection{Model architecture}
\label{app:figures:arch}
\suppfig[0.92\linewidth]{figures/main/Fig_architecture.pdf}
\appfigdesc{Conditional Gaussian mixture variational autoencoder (GMVAE) architecture.
Encoder: spectral EEG/EMG maps pass through CNN ($K{=}[7,5,5,3]$) and MLP blocks. Subject ID feeds a subject embedder used only in the decoder.
Latent: diagonal Gaussian posterior with dimension $d{=}6$ feeds a GMM (cGMVAE) or HMM--GMM (cHMM--GMVAE) prior.
Decoder: latent vector concatenated with subject embedding, then CNN reconstruction. Training minimises mean squared error (MSE) reconstruction plus $\beta \times$ prior KL.
At holdout, only the encoder and prior are used for staging.}
\label{fig_arch}

\subsection{S2 Fig.}
\label{S2_Fig}
\suppfig[0.48\linewidth]{figures/supplementary/S4_per_mouse_nmi.pdf}
\appfigdesc{Per-mouse prior NMI on fold~4. Paired bars for cGMVAE and cHMM-GMVAE on each held-out mouse.}

\subsection{S3 Fig.}
\label{S3_Fig}
\suppfig[0.44\linewidth]{figures/supplementary/S3_confusion.pdf}
\appfigdesc{Macro confusion (fold~4, $K{=}3$, cHMM-GMVAE, best seed).
Row-normalised $3 \times 3$ matrix: expert Wake/NREM/REM (rows) vs.\ predicted macro after Hungarian alignment (columns). REM row shows low recall (${\sim}3\%$) on this holdout fold.}

\subsection{Substage transitions at $K{=}6$}
\suppfig[0.72\linewidth]{figures/main/Fig5_K6_transition.pdf}
\appfigdesc{Substage dynamics at $K{=}6$ (population sweep, best seed).
Change-only transition matrix $P(\mathrm{next} \mid \mathrm{current})$ between mixture components.
Stable within-macro micro-cycles in Wake S2 and NREM S5/S6; short-lived boundary substates S1 and S4 mediate Wake--NREM switches.}
\label{fig5}

\subsection{Population substage PCA ($K{=}4$, $7$, $8$)}
\label{app:population_k}

\noindent Locked cHMM--GMVAE, all 20 mice (incohort), best seed per $K$.
Each panel shows PCA of encoder means with expert Wake/NREM/REM versus predicted substates.

\subsection{S17 Fig.}
\label{S17_Fig}
\suppfig[0.75\linewidth]{figures/supplementary/population/S_pop_K04_pca.png}
\appfigdesc{Population $K{=}4$ (best seed, prior NMI 0.517).
Latent PCA: expert Wake/NREM/REM versus predicted substates.}

\subsection{S18 Fig.}
\label{S18_Fig}
\suppfig[0.75\linewidth]{figures/supplementary/population/S_pop_K07_pca.png}
\appfigdesc{Population $K{=}7$ (best seed, prior NMI 0.468).
Latent PCA: expert macros versus predicted substates.}

\subsection{S19 Fig.}
\label{S19_Fig}
\suppfig[0.75\linewidth]{figures/supplementary/population/S_pop_K08_pca.png}
\appfigdesc{Population $K{=}8$ (best seed, prior NMI 0.439).
Latent PCA: expert macros versus predicted substates.}

\section{Supplementary tables}
\label{app:tables}

\subsection{S10 Table.}
\label{S10_Table}
\apptabdesc{Per-class classification metrics (fold~4, $K{=}3$, cHMM-GMVAE, best seed). Cohen's $\kappa$, macro F1, and per-class recall after Hungarian alignment.}
\apptabular{% Auto-generated by scripts/paper/compute_holdout_classification_metrics.py
\begin{tabular}{|l|r|r|r|r|r|}
\hline
\textbf{Run} & \textbf{$\kappa$} & \textbf{Macro F1} & \textbf{Wake rec.} & \textbf{NREM rec.} & \textbf{REM rec.} \\
\thickhline
fold4 cHMM K=3 & 0.811 & 0.638 & 0.959 & 0.955 & 0.034 \\
\hline
\end{tabular}
}

\subsection{S11 Table.}
\label{S11_Table}
\apptabdesc{Seed stability (joint holdout). Prior NMI mean, standard deviation, and best-of-three per fold and model.}
\apptabular{% Auto-generated by scripts/paper/summarize_holdout_seeds.py
\begin{tabular}{|l|l|r|r|r|r|}
\hline
\textbf{Fold} & \textbf{Model} & \textbf{Seeds} & \textbf{NMI mean} & \textbf{NMI SD} & \textbf{NMI best} \\
\thickhline
1 & cGMVAE & 2 & 0.219 & 0.048 & 0.267 \\
1 & cHMMGMVAE & 3 & 0.521 & 0.041 & 0.579 \\
2 & cGMVAE & 2 & 0.168 & 0.089 & 0.256 \\
2 & cHMMGMVAE & 3 & 0.375 & 0.014 & 0.396 \\
3 & cGMVAE & 2 & 0.297 & 0.009 & 0.307 \\
3 & cHMMGMVAE & 3 & 0.450 & 0.315 & 0.675 \\
4 & cGMVAE & 2 & 0.214 & 0.085 & 0.299 \\
4 & cHMMGMVAE & 3 & 0.581 & 0.056 & 0.649 \\
\hline
\end{tabular}
}

\subsection{S12 Table.}
\label{S12_Table}
\apptabdesc{Substage validity across mixture resolution (population sweep, best seed per $K$).
Active: mixture components with $\geq 1$ assigned epoch.
Min occ.: occupancy of the rarest active substage.
Latent ED: mean pairwise energy distance between substates in encoder space.
Fisher: Fisher trace of latent between-state scatter.
Cram\'{e}r $V$: association between predicted substages and expert macros (post-hoc).
$p_{\mathrm{perm}}$: permutation test for latent separation (shuffle substage labels; $H_0$: observed ED is label-random).
Low $p_{\mathrm{perm}}$ supports distinct substates even when macro NMI falls at high $K$.}
\apptabular{% Auto-generated by scripts/paper/compute_substage_sweep_statistics.py (population)
\begin{tabular}{|r|r|r|r|r|r|r|r|}
\hline
\textbf{$K$} & \textbf{Active} & \textbf{Min occ.\ (\%)} & \textbf{Latent ED} & \textbf{Fisher} & \textbf{Cramér $V$} & \textbf{NMI} & \textbf{$p_{\mathrm{perm}}$} \\
\thickhline
4 & 4/4 & 0.37 & 1.91 & 2.14 & 0.503 & 0.542 & 0.010 \\
5 & 5/5 & 6.19 & 2.61 & 5.98 & 0.591 & 0.484 & 0.010 \\
6 & 6/6 & 3.13 & 2.46 & 7.19 & 0.609 & 0.457 & 0.010 \\
7 & 7/7 & 3.19 & 2.50 & 9.10 & 0.631 & 0.468 & 0.010 \\
8 & 8/8 & 1.38 & 2.20 & 9.62 & 0.659 & 0.439 & 0.010 \\
9 & 9/9 & 1.35 & 2.68 & 11.36 & 0.675 & 0.442 & 0.010 \\
12 & 12/12 & 1.38 & 3.03 & 11.68 & 0.662 & 0.403 & 0.010 \\
15 & 10/15 & 0.01 & 3.48 & 5.16 & 0.496 & 0.409 & 0.010 \\
\hline
\end{tabular}
}

\subsection{S14 Table.}
\label{S14_Table}
\apptabdesc{Holdout prior log-likelihood (joint scope, $K{=}6$, best-of-three seed).
cGMVAE uses sequence length $T{=}1$: each 4\,s epoch is scored independently with the static GMM prior (mean $\log p(z_t)$ per epoch; \texttt{Likelihood} in run metrics).
HMMGMVAE and cHMM--GMVAE use $T{=}64$ contiguous epochs per window; the reported value is the HMM forward marginal mean $\log p(z_{1:T})$ per sequence (\texttt{log p(z\_1:T)}).
Raw totals are \emph{not} comparable across $T{=}1$ and $T{=}64$ models because sequence models accumulate density over 64 steps.
The normalised column divides cGMVAE by $L{=}6$ and sequence models by $T\cdot L{=}384$, matching the orange axis in Fig~\ref{fig2} for $T{=}64$ runs.
Higher values indicate better prior density on held-out encoder means; they need not track NMI when mixture structure misaligns with expert macros.
Fold~4 cGMVAE (seed~2) shows an extreme negative raw value despite moderate NMI---a known case where prior density and clustering quality diverge.}
\apptabular{% Auto-generated by scripts/paper/build_holdout_likelihood_table.py
\begin{tabular}{|l|l|r|r|r|r|}
\hline
\textbf{Fold} & \textbf{Model} & \textbf{$T$} & \textbf{Holdout NMI} & \textbf{Prior LL (raw)} & \textbf{Norm.} \\
\thickhline
1 & cGMVAE & 1 & 0.267 & -17.59 & -2.931 \\
1 & HMMGMVAE & 64 & 0.409 & -3755.5 & -9.780 \\
1 & cHMM--GMVAE & 64 & 0.579 & -26.8 & -0.070 \\
2 & cGMVAE & 1 & 0.256 & -11.86 & -1.976 \\
2 & HMMGMVAE & 64 & 0.236 & -1200.4 & -3.126 \\
2 & cHMM--GMVAE & 64 & 0.396 & -307.3 & -0.800 \\
3 & cGMVAE & 1 & 0.307 & -12.51 & -2.085 \\
3 & HMMGMVAE & 64 & 0.497 & -166.0 & -0.432 \\
3 & cHMM--GMVAE & 64 & 0.675 & -84.5 & -0.220 \\
4 & cGMVAE & 1 & 0.299 & -68904.82 & -11484.137 \\
4 & HMMGMVAE & 64 & 0.404 & -264.6 & -0.689 \\
4 & cHMM--GMVAE & 64 & 0.649 & 281.1 & 0.732 \\
\hline
\end{tabular}
}

\subsection{S15 Table.}
\label{S15_Table}
\apptabdesc{Exploratory substage map ($K{=}12$, population sweep, seed~2).
Mixture index S1--S12 with occupancy, dominant expert macro, biological reading, and proposed merges after sleep-biology review (A.\ Wollesen Jürgensen-Breum, B.\ Kornum).
Dominant macro is post-hoc only; training is fully unsupervised.}
\apptabular{% Biological reading at K=12 (population sweep, seed 2); mixture index S1--S12.
\begin{tabular}{|l|r|l|p{5.8cm}|}
\hline
\textbf{State} & \textbf{Occ. (\%)} & \textbf{Dominant macro} & \textbf{Biological reading and proposed action} \\
\thickhline
S1 & 5.3 & Mixed W/N/R & Short-bout transition bucket; merge with S6 or drop \\
S2 & 3.2 & REM (93\%) & Consolidated REM; model does not split REM further \\
S3 & 25.7 & Wake (99\%) & Core quiet Wake (largest Wake cluster) \\
S4 & 9.2 & Wake (99\%) & Spectrally deep/slow Wake neighbour; review vs.\ N3-like label \\
S5 & 1.4 & NREM (82\%) & High EMG; 5 mice only---artefact / line noise \\
S6 & 2.3 & NREM (78\%) & Light NREM (N1-like); merge with S10 \\
S7 & 2.0 & Wake (98\%) & REM$\to$Wake transition (noisy post-REM epochs) \\
S8 & 7.6 & Wake (90\%) & Short-bout Wake--NREM mix; merge with neighbours \\
S9 & 4.9 & Mixed W/N & Wake--NREM boundary; merge with S6/S10 \\
S10 & 12.6 & NREM (94\%) & Core NREM (N2-like) \\
S11 & 11.5 & NREM (95\%) & High-variability tail of S5-like EMG noise or over-split active Wake \\
S12 & 14.3 & NREM (94\%) & Redundant NREM bucket; collapse into S10 \\
\hline
\end{tabular}
}

\subsection{Supplementary experimental settings}
\label{S1_File}

Full hyperparameters for joint and within-laboratory holdout models are provided as supplementary material.

\end{document}
